\chapter{刑罚的执行和消灭}

% Source PDF physical page 207; see audit/source_coverage.json.
\begin{EditionOpeningPanel}
重要考点：失败的假释与失败的缓刑的对比，追诉时效。
\end{EditionOpeningPanel}

\section{减刑}

第78条被判处管制、拘役、有期徒刑、无期徒刑的犯罪分子，在执行期间，如果认真遵守监规，接受教育改造，确有悔改表现的，或者有立功表现的，可以减刑…..

\subsection*{一、减刑条件}

1.对象条件：被判处\CJKunderwave{管制、拘役、有期徒刑、无期徒刑}的犯罪分子。

（1）至于犯罪行为是故意犯罪还是过失犯罪，是重罪还是轻罪，是危害国家安全罪还是其他犯罪，在所不问。

（2）对缓刑能否减刑？司法解释\footnote{2017年1月1日最高人民法院《关于办理减刑、假释案件具体应用法律的规定》第18条。}规定：对缓刑，一般不适用减刑；若有重大立功，可以减刑，同时缩短缓刑考验期。（2013年第57题）

（3）第78条的减刑针对的是自由刑（包括限制自由刑也即管制），将自由刑的刑期减少。死缓减为无期徒刑或有期徒刑，不属于第78条的减刑，因为第一，死缓不属于自由刑；第二，死缓变成无期徒刑或有期徒刑，属于刑种变更。

2.实际执行刑期的限度条件。刑法没有限制减刑的次数，只要符合减刑条件，可以多次减刑。但是，规定了实际执行刑期的底线。对此问题进行大总结（包括死缓减为无期或有期）。

% Source PDF physical page 208; see audit/source_coverage.json.
\InsertDiagram{figures/manual_table_p208_1.tex}
\SourceFootnote{\textbf{表注①}\quad 多个有期徒刑数罪并罚，最高是25年。25年的1/2，是12.5年。所以，判处无期徒刑的，实际执行刑期不能少于13年，要比有期徒刑多半年。}
\SourceFootnote{\textbf{表注②}\quad 根据刑法第50条第2款规定，对两种死缓犯，可以限制减刑。第一种是累犯。第二种是实施故意杀人、强奸、抢劫、绑架、放火、爆炸、投放危险物质或者有组织的暴力性犯罪的犯罪分子。此外，对累犯和被判10年以上有期徒刑的第二种犯罪分子，不得假释。}

\subsection*{二、减刑程序}

\InsertDiagram{figures/manual_table_p208_2.tex}

\section{假释}

[问题］狗蛋犯诈骗罪，被判有期徒刑5年，罚金5万元。狗蛋有能力缴纳罚金，拒不缴纳。对狗蛋能否适用假释？（不缴纳罚金案）

假释，是指对于被判处有期徒刑、无期徒刑的部分犯罪人，在执行一定刑罚之后，确有悔改表现，不致再危害社会，附条件地予以提前释放的制度。附条件，是指被假释的犯罪人，如果遵守一定条件，就认为原判刑罚已经执行完毕；如果没有遵守一定条件，就收监执行剩余刑罚。

\subsection*{一、假释的条件}

1．对象条件：被判处无期徒刑、有期徒刑的部分犯罪分子。对被判处管制、拘役的，不能假释。

2.禁止适用的两种人：第一种人是累犯；第二种人是实施故意杀人、强奸、抢劫、绑架、放火、爆炸、投放危险物质或者有组织的暴力性犯罪的犯罪分子，并且被判处10年以上有期徒刑、无期徒刑。8种犯罪：烧杀淫掠，绑走暴徒（绑组爆投）。（2014年第55题，2017年第11题）

（1）第一种人和第二种人，被判处死缓，后减为无期徒刑、有期徒刑，不得假释。对这两种死缓犯也可以限制减刑。

（2）第一种人和第二种人之外的罪犯，被判处死缓，后减为无期徒刑、有期徒刑，可

% Source PDF physical page 209; see audit/source_coverage.json.
以假释。例如，甲因贩卖毒品罪被判处死缓，后减为无期徒刑，只要符合假释其他条件，可以假释。

（3）数罪并罚的情形：

\InsertDiagram{figures/manual_table_p209_1.tex}

3.实质条件：确有悔改表现，不再有人身危险性和再犯可能性。

（1）对于财产刑，罪犯确有履行能力而不履行的，不认定其有悔改表现，不予假释。

例如，本节开头“不缴纳罚金案”中，对狗蛋不能适用假释。如果罪犯确无履行能力的，不影响对其确有悔改表现的认定。

（2）对\CJKunderwave{职务犯罪}、破坏金融管理秩序和金融诈骗犯罪、组织（领导、参加、包庇、纵容）黑社会性质组织犯罪等罪犯，\CJKunderwave{不积极退赃}、协助追缴赃款赃物、赔偿损失，不认定其有悔改表现，不予假释。该规定意思是，这三类犯罪“职务犯罪”“金融犯罪”“涉黑犯罪”的罪犯，在退赃等方面必须有积极态度，才可能获得假释。而普通犯罪的罪犯，如果客观上没有履行能力，即使态度不积极，也可能获得假释。（2017年第11题）

（3）罪犯在刑罚执行期间的\CJKunderwave{申诉权利}应当依法保护，对其\CJKunderwave{正当申诉}不能不加分析地认为是不认罪悔罪。

4.在假释考验期必须实行社区矫正。

\subsection*{二、假释的已执行刑期条件和考验期期限}

\InsertDiagram{figures/manual_table_p209_2.tex}

1.例外规定：如果有特殊情况，\CJKunderwave{经最高人民法院}核准，可以不受上述执行刑期的限制。

2.被假释的罪犯，一般不得减刑，也即不得缩短假释考验期。

% Source PDF physical page 210; see audit/source_coverage.json.
\subsection*{三、假释的法律后果}

成功的假释：认为原判刑罚已经执行完毕。

失败的假释：假释被撤销。三种情形以及处理方式：（2015年第12题）

1.发现漏罪。这是指在假释考验期内发现判决宣告以前未判决的罪。

（1）发现的漏罪，是故意犯罪，还是过失犯罪，在所不问；是同种罪，还是异种罪，也在所不问。

（2）撤销假释后，对前罪和漏罪先并后减，也即，将前罪刑期和漏罪刑期进行合并，然后减去前罪已执行刑期。\CJKunderwave{注意}：已经经过的考验期，不算已经执行的刑期。

（3）此后，符合假释其他条件的，可以再次假释。

（4）如果在考验期期满后发现漏罪，就不能撤销假释，只能另行起诉审判。

2.又犯新罪。这是指在考验期内犯新罪。

（1）又犯的新罪，是故意犯罪，还是过失犯罪，在所不问；是同种罪，还是异种罪，也在所不问。

（2）在考验期内犯新罪，无论是在考验期内还是在考验期满后被发现，都要撤销假释。（2013年第57题）

（3）撤销假释后，对前罪和新罪先减后并，也即，先用原判刑期减去已执行刑期，然后用剩余刑期跟新罪刑期合并。\CJKunderwave{注意}：已经经过的考验期，不算已经执行的刑期。

（4）此后，一般不得再次假释。这是因为，又犯新罪，表明仍具有人身危险性和再犯可能性，不符合假释的实质条件。

[提示」如果在缓刑考验期内犯新罪，撤销缓刑，数罪并罚后，有可能适用假释。（2009年第12题）。

[总结】再次适用缓刑、假释问题。

\InsertDiagram{figures/manual_table_p210_1.tex}

（5）在考验期满后犯新罪，不能撤销假释。\CJKunderwave{注意}：可以成立累犯，因为假释考验期满就视为刑罚已经执行完毕。

3.在假释考验期内，违反法律、法规或者监管规定，应当撤销假释，收监执行尚未执行完毕的刑罚。\footnote{（1）在考验期内违反规定，无论是在考验期内还是在考验期满后被发现，都要撤销假释。这一点，与在考验期内犯新罪的处理及理由是相同的。（2）在假释考验期内，没有剥夺犯罪人行使言论、出版、集会、结社、游行、示威自由的权利。这一点和管制不同。执行管制时，未经执行机关批准，不得行使这些权利。假释考验期的遵守规定和缓刑考验期的遵守规定基本相同。（3）此后，一般不得再次假释。}

% Source PDF physical page 211; see audit/source_coverage.json.
\subsection*{四、缓刑和假释的比较}

成功的效果不同：缓刑考验期满，原判刑罚不再执行；假释考验期满，认为原判刑罚已经执行完毕。（2026年试题）

\InsertDiagram{figures/manual_table_p211_1.tex}

\section{刑罚的消灭与追诉时效}

[问题]1996年，南京发生一起杀人碎户案。真凶至今尚未找到。如果2024年找到了，能否追究其刑事责任？（南京碎尸案）

刑罚的消灭，是指由于法定的或事实的原因，致使国家对犯罪人的刑罚权归于消灭。刑罚消灭必须基于一定的事由。\footnote{刑罚消灭事由有：（1）超过追诉时效；（2）经特教令免除刑罚；（3）告诉才处理的犯罪，没有告诉或者撤回告诉；（4）犯罪嫌疑人、被告人死亡。}考试最常考的刑罚消灭事由是追诉时效。

追诉时效，是指刑法规定的，对犯罪人进行刑事追究的有效期限，超过此期限，司法机关就不能追诉。

\subsection*{一、追诉时效的期限}

\InsertDiagram{figures/manual_table_p211_2.tex}

1.这里的法定最高刑，不是指整个法定刑幅度的最高刑，而是犯罪情节对应的量刑幅度（量刑档次、刑格）的最高刑。例如，盗窃罪，数额较大的，处3年以下有期徒刑；数额巨大的，处3年以上10年以下有期徒刑；数额特别巨大的，处10年以上有期徒刑或者无期徒刑。若甲盗窃3000元，则法定最高刑是3年，追诉时效是5年；若甲盗窃500万元，则法定最高刑是无期徒刑，追诉时效是20年。

% Source PDF physical page 212; see audit/source_coverage.json.
2.特别追诉制度，是指如果20年以后，认为必须追诉的，须报请最高人民检察院核准。注意，其他法定刑档次，追诉时效过了，就没有救济机会了。例如，本节开头的“南京碎户案”，超过了20年追诉期限，要追诉的话，须报请最高人民检察院核准。（2023年试题）

3.单位犯罪的追诉期限一律是5年，因为单位犯罪的法定刑只能是罚金，罚金处在“不满5年有期徒刑”这一档。对直接责任人的追诉期限按照自然人的法定刑计算。

4.共同犯罪。（1）共同犯罪人的量刑档次不同的话，追诉期限也不同。例如，实行犯与帮助犯的量刑档次不同，追诉期限也不同。（2024年试题）

（2）对各共犯人分别计算各自的追诉时效。一人超过追诉时效，另一人没有超过，则只能对后者追诉。例如，甲、乙共同杀人，甲应是主犯，乙应是从犯，甲应被判无期徒刑，乙应被判10年。经过15年后，只能追诉甲，不能追诉乙。（2016年第10题，2018年试题)

5.想象竞合犯、牵连犯触犯了两个罪。如果轻罪已过追诉时效，而重罪未过，则只按该重罪追诉。例如，甲先伪造国家机关证件，后持该证件实施诈骗，数额巨大。经过5年后，伪造国家机关证件罪已过追诉时效，诈骗罪未过，则不再追诉前者，只追诉后者。

\subsection*{二、追诉时效的计算}

\InsertDiagram{figures/manual_table_p212_1.tex}

\subsection*{三、追诉时效的延长}

这是指在追诉时效的进行期间，因为发生法律规定的事由，而使追诉时效暂时停止执行。包括两种情形：

1.第88条第1款：“在人民检察院、公安机关、国家安全机关立案侦查或者在人民法院受理案件以后，\CJKunderwave{逃避侦查或者审判的}，不受追诉期限的限制。”（2023年试题）

例如，甲犯某罪，公安机关立案侦查后，甲误以为自已没犯罪，去外地打工。甲不属于逃避侦查。

2.第88条第2款：“被害人在追诉期限内提出控告，人民法院、人民检察院、公安机关\CJKunderwave{应当立案而不予立案的}，不受追诉期限的限制。”

例如（2015年第60题），李某被人打了，过了追诉期限才报案，公安机关未立案。对此，受追诉期限限制，不能追诉。

[注意]上述两种情况虽然不受追诉期限的限制，但如果再犯新罪，新罪仍然受追诉期限的限制（各算各的账）。例如，行为人的甲罪被司法机关立案侦查，但行为人逃避侦查与审判，其后又犯了乙罪。先前的甲罪虽然不受追诉期限的限制，但后来的乙罪仍受追诉期限的限制。

% Source PDF physical page 213; see audit/source_coverage.json.
\subsection*{四、追诉时效的中断}

第89条第2款在\CJKunderwave{追诉期限以内}又犯罪的，\CJKunderwave{前罪追诉的期限从犯后罪之日起计算}。

追诉时效的中断，也称为追诉时效的重新起算，是指在追诉时效进行期间，犯罪人又犯罪，导致以前所经过的时效期间归于无效，追诉时效重新起算。（2015年第60题，2016年第10题，2018年试题）

例如，甲犯A罪，追诉时效是15年。在第8年，甲又犯B罪。此时，A罪经过的8年追诉期间作废（清零），A罪的15年追诉期限重新起算。

\InsertDiagram{figures/manual_p213_1.tex}

1．各算各的账。

（1）共同犯罪中，各共犯人的追诉时效的中断，互不影响。例如，甲、乙共同犯罪，在追诉期限内，甲又犯新罪，则甲的前罪的追诉期限重新计算，乙的追诉期限照旧。

（2）牵连犯中，前面手段行为的犯罪的追诉期限会被后面目的行为的犯罪中断而重新计算。

（3）前罪继续向前，后罪已经到期的情形。例如，甲于2000年1月犯普通抢劫罪，追诉期限是15年；甲于2005年1月又犯普通盗窃罪，追诉期限是5年。此时，抢劫罪的追诉期限要从2005年1月重新计算。到了2012年1月，盗窃罪已经超过追诉期限，而抢劫罪还没过，此时只追究抢劫罪。

2.追诉时效的中断与延长并存时，只适用追诉时效延长的规定。例如，甲于2000年1月犯盗窃罪，公安机关立案侦查后逃避侦查，并于2002年1月犯诈骗罪。此时，对盗窃罪不能适用追诉时效中断的规定，而应适用追诉时效延长的规定。

3.空间效力。后罪应是我国刑法规定的犯罪，并有管辖权。例如，甲杀人后逃往国外17年时，在国外伪造私人印章，我国刑法没有将此规定为犯罪，那么该行为就不会造成故意杀人罪的追诉时效的中断（更新）。

4.时间效力。现行刑法与1979年刑法关于追诉时效的条文规定有所不同。两种条文的时间效力，应像一般条文一样，遵循“从旧兼从轻”的原则。（2016年第10题）

[总结]两高的特殊权力：（2012年第13题）

\InsertDiagram{figures/manual_table_p213_2.tex}

% Source PDF physical page 214; see audit/source_coverage.json.
\ContinuedTable{figures/manual_table_p214_1.tex}
